\documentclass[11pt]{amsart}

\usepackage[T1]{fontenc}
\usepackage{lmodern}
\usepackage{microtype}
\usepackage{mathtools,amssymb,amsthm}
\usepackage[hidelinks]{hyperref}

\hypersetup{
  pdftitle={A counterexample, with a hyperbolic summand, to a printed
  subadditivity question for topological volume}
}

\newtheorem{theorem}{Theorem}
\newtheorem{lemma}[theorem]{Lemma}
\theoremstyle{definition}
\newtheorem*{question}{Question}
\theoremstyle{remark}
\newtheorem*{remark}{Remark}

\newcommand{\volt}{\operatorname{vol_t}}
\newcommand{\volh}{\operatorname{vol_h}}
\newcommand{\HH}{\mathbb{H}}

\title[A counterexample to a subadditivity question]
{A counterexample, with a hyperbolic summand, to a printed subadditivity
question for topological volume}

\date{}

\subjclass[2020]{57K32, 57K30, 57M50}
\keywords{topological volume, hyperbolic volume, connected sum, Weeks manifold,
counterexample}

\begin{document}

\begin{abstract}
For a closed orientable $3$--manifold $M$, Kegel, Ray, Spreer, Em Thompson and Tillmann
define the \emph{topological volume} $\volt(M)$ to be the infimum of $\volh(M\setminus L)$
over the hyperbolic links $L\subset M$, the empty link included, and ask whether
$\volt(M\#N)\le\volt(M)+\volt(N)$ for all closed orientable $M,N$. We exhibit one pair for
which this fails, so the question is false exactly as printed. Let $W$ be the Weeks
manifold. Then $W$ is hyperbolic, so $\volt(W)=\volh(W)\le 0.943$, whereas $W\#W$ is
reducible and therefore not hyperbolic, so every competitor in its infimum is an
\emph{orientable cusped} finite-volume hyperbolic $3$--manifold and Cao--Meyerhoff gives
$\volt(W\#W)\ge 2v_0=2.029883212819307\ldots>1.886\ge\volt(W)+\volt(W)$. Both summands here
are hyperbolic; the reading of the question in which the summands are not hyperbolic, which
is the substantive connected-sum problem, is untouched.
\end{abstract}

\maketitle

\section{The question}

Let $M$ be a closed orientable $3$--manifold. A \emph{link} in $M$ is a possibly empty,
possibly disconnected closed $1$--submanifold $L\subset M$, and $L$ is \emph{hyperbolic} if
$M\setminus L$ admits a complete hyperbolic structure of finite volume. Kegel, Ray, Spreer,
Em Thompson and Tillmann \cite{KRSTT} define
\begin{equation}\label{eq:def}
  \volt(M)=\inf\bigl\{\,\volh(M\setminus L)\;\big|\;L\subset M
  \text{ is a hyperbolic link}\,\bigr\},
\end{equation}
following Bessi\`eres, Besson, Boileau, Maillot and Porti \cite{BBBMP}. Because the empty
link is admitted, they record immediately after \eqref{eq:def} that
\begin{equation}\label{eq:hyp}
  \volt(M)=\volh(M)\qquad\text{whenever $M$ is hyperbolic,}
\end{equation}
in their own words: \emph{``If $M$ admits a complete hyperbolic structure, then
$\volt(M)=\volh(M)$, since volume decreases under Dehn filling.''}
Write $v_0$ for the volume of the regular ideal hyperbolic tetrahedron, so that
\begin{align*}
  v_0&=1.014941606409653625021203\ldots,\\
  2v_0&=2.029883212819307250042405\ldots
\end{align*}

In the subsection \emph{Connected sums} of \cite{KRSTT} the authors ask the following,
quoted verbatim.

\begin{question}[\cite{KRSTT}, first Question of the subsection \emph{Connected sums}]
``Let $M$ and $N$ be closed, orientable $3$-manifolds. Is
$\volt(M\# N)\le\volt(M)+\volt(N)$?''
\end{question}

As printed, the Question carries exactly two hypotheses, \emph{closed} and
\emph{orientable}: no primeness, no irreducibility and no hyperbolicity clause. We found no
standing convention elsewhere in \cite{KRSTT} that would add one.

\section{Two lemmas}

\begin{lemma}[a floor at non-hyperbolicity]\label{lem:floor}
Let $X$ be a closed orientable $3$--manifold that is not hyperbolic. Then
$\volt(X)\ge 2v_0$.
\end{lemma}

\begin{proof}
Let $L\subset X$ be a hyperbolic link. Then $L\ne\emptyset$: otherwise $X$ itself would
carry a complete finite-volume hyperbolic structure, contrary to hypothesis. Hence
$X\setminus L$ is a non-compact --- so cusped --- complete finite-volume hyperbolic
$3$--manifold, and it is orientable, orientability being inherited from $X$. By
Cao--Meyerhoff \cite{CM}, \emph{every orientable cusped hyperbolic $3$--manifold has volume
at least $2v_0$}, whence $\volh(X\setminus L)\ge 2v_0$. Take the infimum over $L$.
\end{proof}

\begin{remark}
The hypothesis \emph{orientable} in \cite{CM} is load-bearing: the Gieseking manifold is a
non-orientable cusped hyperbolic $3$--manifold of volume
$v_0=1.014941606409653\ldots<1.886$, so a floor of $v_0$ would not suffice below.
Lemma~\ref{lem:floor}
does not presuppose that the infimum in \eqref{eq:def} is attained, nor that it is finite:
every competitor is bounded below by $2v_0$, so the infimum is, whatever it is.
\end{remark}

\begin{lemma}[reducible $\Rightarrow$ not hyperbolic]\label{lem:red}
If neither $M$ nor $N$ is $S^3$, then $M\#N$ is reducible and is not hyperbolic.
\end{lemma}

\begin{proof}
The summing sphere bounds a ball on neither side, so it is essential and $M\#N$ is
reducible. A closed hyperbolic $3$--manifold has universal cover $\HH^3$, hence is
aspherical, hence has $\pi_2=0$, hence contains no essential embedded $2$--sphere by the
Sphere Theorem: it is irreducible.
\end{proof}

\section{The counterexample}

Let $W$ be the Weeks manifold. It is closed, orientable and hyperbolic, so
$\volt(W)=\volh(W)$ by \eqref{eq:hyp}, and
\begin{equation}\label{eq:weeks}
  \volh(W)\le 0.943
\end{equation}
by Gabai--Meyerhoff--Milley, whose abstract reads, verbatim: \emph{``We also show how this
result can be used to construct a complete list of all one-cusped hyperbolic
three-manifolds with volume $\le 2.848$ and all closed hyperbolic three-manifolds with
volume $\le 0.943$. In particular, the Weeks manifold is the unique smallest volume closed
orientable hyperbolic $3$-manifold.''} We use only the \emph{bound} \eqref{eq:weeks}, never
a decimal expansion of $\volh(W)$.

\begin{theorem}\label{thm:A}
$\volt(W\# W)\ge 2v_0>1.886\ge\volt(W)+\volt(W)$. In particular the Question of \S1 is
false at the pair $(W,W)$.
\end{theorem}

\begin{proof}
By \eqref{eq:hyp} and \eqref{eq:weeks}, $\volt(W)+\volt(W)\le 2\times 0.943=1.886$. Neither
summand is $S^3$, so $W\# W$ is not hyperbolic by Lemma~\ref{lem:red}, and
Lemma~\ref{lem:floor} applied to $X=W\# W$ gives $\volt(W\# W)\ge 2v_0$. Finally
$2v_0-1.886\ge 0.143883212819307>0$.
\end{proof}

\section{What is not settled}

Both summands of the counterexample are hyperbolic, and that is what the argument uses: the
substantive connected-sum problem, in which the summands are not hyperbolic, is untouched.
There Lemma~\ref{lem:floor} already gives $\volt(M),\volt(N)\ge 2v_0$, so a counterexample
would need a lower bound on $\volt(M\# N)$ exceeding a sum of at least $4v_0$, and nothing
above supplies one. The authors of \cite{KRSTT} state \eqref{eq:hyp} and quote
Cao--Meyerhoff's $2v_0$ minimum themselves, so a hypothesis of primeness, or of
non-hyperbolic summands, may well have been intended and omitted; we settle the question as
printed and make no claim about what was meant. Nothing above supports or refutes
$\volt(M\#N)\le\volt(M)+\volt(N)+2v_0$ or any other repaired form. We exhibit the one pair
$(W,W)$ and count no others. The two mathematical inputs beyond \eqref{eq:weeks} are the
source's own: \eqref{eq:hyp} is its sentence, and it quotes Cao--Meyerhoff itself.

\medskip
\noindent\textbf{Adjacent work.} The subsection \emph{Connected sums} of \cite{KRSTT}
records
\[
  \volt(L(2,1)\# L(3,1))<\volt(L(2,1))+\volt(L(3,1)),
\]
an instance in which the
inequality of the Question of \S1 holds strictly, and the summary of related work in
\cite{KRSTT} says on the strength of it that $\volt$ ``is not additive under connected
sum''; that sentence is about that strict inequality and not about any failure of the
inequality it asks about, which \cite{KRSTT} leaves open. Example 2.1 of
\cite{KRSTT} computes $\volt(S^3)=2v_0$ from the same minimum-volume input, which is the
argument of Lemma~\ref{lem:floor} run for the single manifold $S^3$; we did not find the
general statement of Lemma~\ref{lem:floor} there. The
Oberwolfach report \cite{Oberwolfach}, by a working group including a co-author of
\cite{KRSTT}, closes by asking to better understand the behaviour of hyperbolic volume under
connected sums, and settles nothing about the inequality above.

\section{Verification}

The result above is a hand proof: two lemmas from published theorems, then arithmetic on
three decimal places. Nothing in \S\S1--3 needs a computer.

A program (\texttt{verify.py}, with its transcript \texttt{verify.output.txt}) accompanies
this note because the note prints numbers. It is standard-library Python, and in exact
arithmetic it recomputes $v_0$ and $2v_0$ from $2v_0=3\operatorname{Cl}_2(2\pi/3)$ by the
Clausen--Bernoulli expansion with an explicit truncation bound, cross-checked against
$v_0=\operatorname{Cl}_2(\pi/3)$ via the Clausen duplication identity, and re-decides
$2v_0>1.886$ and $v_0<1.886$ as comparisons of exact rationals. Those digit expansions are
the only values in this note that come from computation rather than from print; the bound
$0.943$ is transcribed, and no decimal expansion of $\volh(W)$ is used anywhere. The program
does not, and cannot, check the topological content quoted from the literature ---
\eqref{eq:hyp}, Cao--Meyerhoff, Gabai--Meyerhoff--Milley, the Sphere Theorem and
Kneser--Milnor are imported, and its own transcript says so.

The transcript's $38$ checks also cover material this note does not state, and they are not
evidence for anything it does state: rows of Table~1 of \cite{KRSTT}; a second route through
a census reach of $3.07$ quoted from a theorem this note does not use, together with ten
further pairs $W\# N$ built on that route, whose premise --- irreducibility of those rows ---
the run leaves unverified; by-product ratios; and one floating-point SnapPy measurement that
the run itself flags as not interval-certified. Theorem~\ref{thm:A} uses none of them, and
this note asserts none of them.

\begin{thebibliography}{9}

\bibitem{BBBMP}
L.~Bessi\`eres, G.~Besson, M.~Boileau, S.~Maillot, and J.~Porti,
\emph{Geometrisation of $3$-manifolds},
EMS Tracts in Mathematics \textbf{13}, European Mathematical Society, 2010,
\S1.3.2. \href{https://doi.org/10.4171/082}{doi:10.4171/082}.
Not open access; consulted only through works that restate its definition of the invariant,
including \cite{KRSTT}.

\bibitem{CM}
C.~Cao and G.~R. Meyerhoff,
\emph{The orientable cusped hyperbolic $3$-manifolds of minimum volume},
Invent. Math. \textbf{146} (2001), no.~3, 451--478.
\href{https://doi.org/10.1007/s002220100167}{doi:10.1007/s002220100167}; MR1869847.
The hypothesis \emph{orientable} in the title is the one used in Lemma~\ref{lem:floor}.

\bibitem{GMM}
D.~Gabai, R.~Meyerhoff, and P.~Milley,
\emph{Minimum volume cusped hyperbolic three-manifolds},
\href{https://arxiv.org/abs/0705.4325}{arXiv:0705.4325}.
The bound \eqref{eq:weeks} is the sentence of the abstract quoted in \S3, read at that
address; that quotation, and not a journal volume or page, is the citation this note relies
on. We give no journal coordinates for this item because we did not check any.

\bibitem{KRSTT}
M.~Kegel, A.~Ray, J.~Spreer, E.~Thompson, and S.~Tillmann,
\emph{On a volume invariant of $3$-manifolds},
Expo. Math. \textbf{44} (2026), no.~4, Paper No.~125782, 31 pp.
\href{https://doi.org/10.1016/j.exmath.2026.125782}{doi:10.1016/j.exmath.2026.125782};
preprint \href{https://arxiv.org/abs/2402.04839v2}{arXiv:2402.04839v2}.
All quotations in this note are from the \texttt{v2} e-print, the version revised following
a referee report.

\bibitem{Oberwolfach}
K.~Ben Aribi, B.~Burton, S.~Ibarra, J.~Morgan, J.~Purcell, J.~S. Quintanilha, F.~Santoro,
S.~Schleimer, and E.~Thompson,
\emph{Computing the topological volume of some three-manifolds},
Oberwolfach Rep. \textbf{21} (2024), no.~3, 2530--2534.
\href{https://doi.org/10.4171/OWR/2024/43}{doi:10.4171/OWR/2024/43}.
\emph{Caution:} the coordinates \cite{KRSTT} gives for this report do not
resolve; the item is identified here by title and author list.

\end{thebibliography}

\end{document}
